\section{Conclusion}
\label{sec:conclusion}
\Sys\ shows that a useful dual-process assistant fits on four CPU cores: a
small Qwen model answers at once, a gate derived from expected loss decides
when deliberation pays, a reasoning model deliberates under a budget, exact
verifiers decide what counts as known, and an audited ledger keeps the system
honest about itself. The mathematics fixed the design before the
measurements did: precision on checkable tasks is the verifier's, not the
model's; agreement is a flag and never proof; every loop is bounded; and a
loop that accepts only on a verifier's pass is stable whatever the model
does inside it.

Three lessons recur across the measurements. First, on a CPU the gains are
in not repeating work. Decode runs at the memory roofline, so more threads,
speculation and hedging did not pay, while prefix reuse, shared prefill
across candidates, one pass over the weights for several users and an
in-process index did, by factors from 1.2 to three orders of magnitude.
Second, what makes small models usable as agents is discipline imposed by
the harness: an action language the engine enforces, a history that is never
rewritten, an exact verifier after every change and a hard budget. Third,
cheap candidates from the quick model, judged by an exact verifier before the
expensive model is consulted, are the best use of idle cores, in the chat
path and in the agent alike.

The open work is stated in Section~\ref{sec:limits}: a ledger large enough to
promote the calibrated gate, retrieval for reworded questions, evaluation at
release scale and on public benchmarks, and the training steps that need a
GPU. The system, its tests and every stored run behind the figures in this
paper are available with the source.
